针对问题二，模型链条从上游资源-鱼群动态轨迹开始：先据此得到综合生态状态指数，再由该指数驱动中层猎物承载力，最后建立长江江豚、长江鲟及入侵物种种群响应模型。高阻隔、无人工放流和污染被设置为三种对照情景。基线情景下，长江江豚与长江鲟种群末值分别为禁渔前的1.133倍和1.140倍；高阻隔情景下分别降至0.901倍和0.856倍；无人工放流时长江鲟降至0.519倍；污染情景下二者进一步降至0.610倍和0.704倍。结果由此落到两点：江湖通道阻隔会同步抑制两类珍稀物种恢复，人工增殖放流对长江鲟的补偿效应更为显著。